\documentclass[10pt,a4paper]{amsart}
\usepackage[margin=19mm]{geometry}
\usepackage{amsmath,amssymb,mathtools}
\usepackage[scr=rsfs,cal=euler]{mathalfa}
\usepackage{xfrac}
\usepackage[unicode,hidelinks,bookmarks=false]{hyperref}
\newcommand{\ie}{i.e.\ }
\newcommand{\cf}{cf.\ }
\newcommand{\resp}{respectively\ }
\newcommand{\Subsection}{\subsection}
\providecommand{\nobreakdash}{\nobreak\hbox{-}}
\setlength{\emergencystretch}{2em}
\multlinegap0pt
\usepackage{bm}
\usepackage{bbm}
\usepackage{dsfont}
\usepackage{xspace}
\usepackage{cancel}
\usepackage{mathtools}
\usepackage{amsthm}
\makeatletter
\@ifundefined{theorem}{\newtheorem{theorem}{Theorem}}{}
\makeatother
\title[Diameter reduction via arc reversal]{Diameter reduction via arc reversal}
\author{Panna Geh\'er, Max K\"olbl, Lydia Mirabel Mendoza-Cadena, Daniel P. Szabo}
\date{Source: arXiv:2402.06259v2 (CC BY 4.0)}
\begin{document}
\maketitle

\noindent\emph{Source and licence.} Diameter reduction via arc reversal. Version arXiv:2402.06259v2 (CC BY 4.0), distributed under the Creative Commons Attribution 4.0 International licence, as recorded in the arXiv licence metadata for this exact version and shown on the abstract page https://arxiv.org/abs/2402.06259v2. Full rights notice: licenses/arxiv-2402.06259-CC-BY-4.0.txt.\par\smallskip

\noindent\textit{Editorial scope.} Counted unit: the Theorem whose source label is \texttt{None} in the article (source0.tex lines 936--938, source label \texttt{None}), delivered together with the article's own complete original proof of it (source0.tex lines 940--957, one \texttt{proof} environment of 18 lines). No mathematics is written, repaired or re-derived: statement and proof are the article's own text line for line. Required context retained uncounted: none. Every definition, display and label of those lines is retained unchanged. Omitted from this package and not redistributed: 1--935, 939--939, 958--967. Nothing inside a retained excerpt is omitted or altered: every retained body is delivered exactly as the article's lines are written, so no omission receipt of any kind accompanies this package. Citations: the retained excerpts carry no citation command at all, so no bibliography entry of the article is transcribed and no reference is redistributed. Reference adaptation: none was needed and none was made; every reference inside the delivered unit resolves inside it, so no unresolved reference or citation is produced. Printed slips kept literal: no printed word, symbol, hypothesis or number of the counted statement or of its proof was altered. Layout and class: the article's own class is \texttt{article} with packages including enumerate, mathtools, bm, cite, enumitem, caption, amsmath, amssymb, amsthm, graphicx, textcomp, fontenc, marvosym, authblk; the wrapper is the lane's pinned \texttt{amsart} editorial wrapper at 10pt on A4, which restates the theorem-like environments the retained text needs and adds the article's own macro definitions that the retained text needs (none), with extra wrapper packages (bm, bbm, dsfont, xspace, cancel, mathtools). The delivered numbering is the unit's own. Assets: no retained span contains a figure environment, a graphics-inclusion command, a PDF-page inclusion, a table or a TikZ picture; no image file is copied into this package, the package declares no asset dependency and no image file is redistributed with it. Licence: version arXiv:2402.06259v2 of this article is distributed under the Creative Commons Attribution 4.0 International licence, as recorded in the arXiv licence metadata for this exact version and shown on the abstract page https://arxiv.org/abs/2402.06259v2. A full rights notice is delivered at licenses/arxiv-2402.06259-CC-BY-4.0.txt. Characters: every retained excerpt is pure ASCII, so the delivered engine, XeLaTeX with its default Latin Modern text font, reports no missing character.
\begin{theorem}
    For every integer $i\geq 8$ there exists an integer $j<i$ and a pair of directed graphs $G, \, H$ with diameter $i$ and $j$ respectively such that their respective directed edge polytopes have the same volume.
\end{theorem}

\begin{proof}
    We start by noticing that directed cycles correspond to standard reflexive simplices and that reversing the direction corresponds to a unimodular (and hence in particular volume preserving) transformation of the associated simplex.
    Further, we notice that joining two graphs at a vertex corresponds to taking the free sum of the corresponding simplices, meaning that the volumes of their corresponding directed edge polytopes are multiplied.
    Hence, the volume of the directed edge polytope of a directed cactus graph does not depend on the direction of the individual cycles.

    Now we can construct $G$ and $H$.
    First assume $i$ was even.
    Write $i=2k$.
    For $G$, consider a directed path of length $k$.
    To every edge in that path, attach a path of length $2$ parallel to its ends and direct it the other way so that it now lies in a directed $3$-cycle.
    This results in a directed cactus graph of $3$-cycles with diameter $2k$.
    Reversing exactly one of the cycles that is attached to two other cycles results in the graph $H$ which has diameter $(2k-1)$.

    If $i$ is odd, we can use an analogous construction.
    Write $i=2k+1$.
    For $G$, we start with a directed path of length $k$ again and complete every edge to a $3$-cycle, except the last one which shall be completed to a $4$-cycle instead.
    The construction of $H$ is analogous.
\end{proof}

\end{document}
